\documentclass[../main.tex]{subfiles}
\begin{document}

\section{Disposiciones generales}

\subsection{Documentación}

Toda la documentación adjunta referente al proyecto comprende los siguientes documentos:

Memoria, Anexos, Planos, Pliego de condiciones y Presupuesto.

\subsection{Normas y Referencias aplicables}

Las normas y referencias aplicables se especifican en el documento de la memoria.

\subsection{Refrigerantes}

Los refrigerantes empleados en el proyecto se especifican en el documento de la memoria.

\subsection{Condiciones para las soldaduras}

Todas las soldaduras que se realicen durante el proyecto deben ejecutarse con la máquina térmica libre de refrigerante, dado que estos presentan una alta inflamabilidad que compromete la seguridad de la operación.

\newpage

\section{Carga de refrigerante en la máquina}

La carga del refrigerante se realizará siempre por el punto de aspiración del compresor, ya que, por razones de seguridad, dicha carga debe efectuarse en un punto situado en baja presión. En los ensayos experimentales se incorpora una tubería adicional que permite la entrada del refrigerante al circuito.

La carga de los refrigerantes debe realizarse siempre desde los envases reglamentarios destinados al transporte y almacenamiento de los refrigerantes correspondientes. Antes de iniciar el proceso de carga, es imprescindible evacuar el refrigerante en uso mediante una bomba de vacío habilitada para tal fin.

Para el proceso de carga de refrigerante se empleará una báscula de precisión con una incertidumbre de medida máxima de 1 gramo, sobre la que se colocará la botella del refrigerante para comprobar la cantidad de masa introducida en la instalación.

Una vez finalizada la operación de carga, bajo ninguna circunstancia podrá permanecer la botella conectada a la máquina durante el funcionamiento de esta.

\newpage

\section{Fluidos refrigerantes}

\subsection{Propano (R290)}

El propano (R290) es un refrigerante natural perteneciente a la familia de los hidrocarburos (alcanos). Es altamente inflamable (categoría ASHRAE A3), por lo que su uso está limitado a instalaciones de pequeñas dimensiones con baja carga de refrigerante. Este gas deberá mantenerse en todo momento alejado de herramientas, procesos con altas temperaturas, llamas o posibles fuentes de ignición, y se realizará un control continuo de los depósitos de propano para prevenir cualquier fuga.

\subsection{Propileno (R1270)}

El propileno (R1270) es un refrigerante natural perteneciente a la familia de los hidrocarburos (alquenos). Es altamente inflamable (categoría ASHRAE A3), por lo que su uso está limitado a instalaciones de pequeñas dimensiones con baja carga de refrigerante. Este gas deberá mantenerse en todo momento alejado de herramientas, procesos con altas temperaturas, llamas o posibles fuentes de ignición, y se realizará un control continuo de los depósitos de propileno para prevenir cualquier fuga.

\subsection{Dimetil éter (RE170)}

El dimetil éter (RE170) es un refrigerante natural perteneciente a la familia de los éteres. Es altamente inflamable (categoría ASHRAE A3), por lo que su uso está limitado a instalaciones de pequeñas dimensiones con baja carga de refrigerante. Este gas deberá mantenerse en todo momento alejado de herramientas, procesos con altas temperaturas, llamas o posibles fuentes de ignición, y se realizará un control continuo de los depósitos de dimetil éter para prevenir cualquier fuga.

\subsection{Especificaciones de seguridad de los refrigerantes utilizados}

Los refrigerantes empleados durante el desarrollo de los ensayos son el R290, R1270 y RE170. El etiquetado de estos refrigerantes es similar e incluye las siguientes advertencias:

\begin{itemize}
    \item Gas extremadamente inflamable.
    \item Contenido a presión: puede explotar en caso de calentamiento.
\end{itemize}

\begin{figure}[ht]
    \centering
    \includegraphics[width=0.65\textwidth]{imagenes/etiquetas_seguridad.png}
    \caption{Etiquetas de peligros}
    \label{fig:etiquetas_seguridad}
\end{figure}

En consecuencia, en el presente apartado del pliego de condiciones se establecen las especificaciones de seguridad mínimas que deberán observarse durante el desarrollo del proyecto, prestando especial atención a la normativa para atmósferas explosivas ATEX.

\begin{itemize}
    \item Verificar la ausencia de fugas en el circuito que puedan provocar ignición. En los cambios de refrigerante se realizarán 30 minutos de vacío mediante una bomba de vacío, lo que permite extraer el refrigerante impregnado en el aceite. Dicha bomba deberá ser adecuada para uso con gases inflamables —por ejemplo, sin motor de escobillas— o estar ubicada de forma que no quede expuesta a las emisiones de refrigerantes inflamables.
    \item No introducir ninguna fuente potencial de ignición que pueda encender un escape de refrigerante.
    \item Utilizar sistemas electrónicos de detección de gases específicamente diseñados para uso con gases inflamables.
    \item Disponer de sistemas de ventilación que contribuyan a dispersar cualquier escape de refrigerante.
    \item Señalizar adecuadamente mediante etiquetas, de modo que todo el personal técnico del laboratorio o taller sea consciente de que la máquina opera con refrigerante inflamable.
    \item Emplear una balanza para la carga de refrigerante que, en caso de ser electrónica, deberá estar homologada para uso en zonas donde puedan estar presentes refrigerantes inflamables, según confirme el fabricante.
    \item Usar una manguera de ventilación para el vaciado de los refrigerantes inflamables, especialmente en el caso de los hidrocarburos.
    \item Disponer de un extintor de incendios a menos de 5 metros de la ubicación del prototipo.
    \item El personal encargado de la carga de refrigerantes deberá llevar guantes y gafas protectoras durante todo el proceso.
    \item Controlar la temperatura de la zona de almacenamiento de las botellas de refrigerantes homologadas, para evitar que se calienten, dado el riesgo potencial de explosión.
    \item Las botellas deberán almacenarse en una zona sin contacto con fuentes de calor ni luz solar directa, en la que esté prohibido fumar y que cuente con ventilación suficiente.
    \item En caso de contacto con la piel: lavar la zona afectada con abundante agua. Retirar la ropa y el calzado contaminados. Para minimizar el riesgo de descargas estáticas e ignición del gas, humedecer completamente la ropa contaminada antes de retirarla. Buscar atención médica si aparecen síntomas. Lavar la ropa antes de volver a utilizarla y limpiar el calzado a fondo antes de reutilizarlo.
    \item En caso de contacto directo con los ojos: enjuagar inmediatamente con abundante agua, levantando ocasionalmente los párpados superior e inferior. Verificar si la víctima usa lentes de contacto y, en caso afirmativo, retirarlos. Continuar el enjuague durante al menos 10 minutos y obtener atención médica si se produce irritación.
    \item En caso de inhalación: trasladar a la víctima al aire libre y mantenerla en reposo en una posición que facilite la respiración. Si se sospecha que los vapores continúan presentes, el personal de rescate deberá usar una mascarilla adecuada o un equipo de respiración autónomo. Si la respiración es inexistente, irregular o se produce un paro respiratorio, el personal capacitado deberá administrar respiración artificial u oxígeno. La respiración boca a boca puede resultar peligrosa para el socorrista. Buscar atención médica si las condiciones adversas persisten o son graves. Si la víctima está inconsciente, colocarla en posición lateral de seguridad y solicitar atención médica de inmediato. Garantizar una buena circulación de aire y aflojar cualquier prenda que oprima, como el cuello de la camisa, corbata o cinturón.
\end{itemize}

\newpage

\section{Materiales y componentes}

\subsection{Compresor}

El compresor deberá ser el referenciado en la memoria, o un equivalente en buen estado y de tipo hermético, acompañado del variador de frecuencia correspondiente.

Para su regulación y protección, dispondrá de un presostato de alta y uno de baja presión. Si la normativa aplicable así lo exige, se instalará también una válvula de seguridad en la descarga. El separador de aceite está incorporado en los compresores seleccionados; en caso contrario, será necesario añadir dicho separador junto con el acumulador de aceite.

La instalación se dispondrá de forma que todas sus partes sean accesibles al personal técnico para las tareas de reparación o mantenimiento que pudieran requerirse.

\subsection{Condensador}

El condensador será el seleccionado en la memoria o, en caso de emplear un equivalente, este deberá contar con la capacidad térmica suficiente para atender las necesidades de la instalación.

\subsection{Evaporador}

El evaporador deberá disipar la mayor cantidad de calor posible; por ello, esta parte del circuito no estará aislada térmicamente. Dado que las temperaturas de trabajo son inferiores a 0 °C, serán necesarios sistemas de desescarche.

\subsection{Válvula de expansión}

La válvula de expansión deberá ser la seleccionada en la memoria o, en caso de emplear una equivalente, esta deberá ser capaz de adaptarse a los diferentes rangos de pérdida de carga requeridos por la instalación.

\subsection{Tuberías}

Las tuberías serán de cobre pulido y se evitará la presencia de vapor de agua en su interior. Las uniones se realizarán mediante soldadura oxiacetilénica, comúnmente conocida como soldadura fuerte. En las uniones de cobre con otros materiales, se empleará decapante para garantizar la correcta ejecución de la soldadura y asegurar la estanqueidad del circuito. En todo el circuito, las tuberías deberán aislarse para evitar pérdidas de capacidad calorífica y la posible formación de hielo en su superficie. Su trazado se dispondrá en zonas resguardadas, evitando áreas de paso.

\newpage

\section{Manipulación y mantenimiento de la planta}

\subsection{Detección y control de fugas}

Los operadores de aparatos que contengan gases deberán adoptar las medidas técnica y económicamente viables para minimizar las fugas a la atmósfera y evitar así la liberación no intencionada de dichos gases. Cuando se detecte una fuga, los operadores garantizarán que los aparatos sean reparados sin demora injustificada.

\subsection{Recuperación del gas}

Cuando se proceda a la extracción del gas de la planta, los operadores deberán asegurarse de que sea recuperado para su reciclaje, regeneración o destrucción. La recuperación deberá llevarse a cabo por una persona física con la certificación requerida.

\subsection{Formación y certificación}

El personal que vaya a manipular la instalación deberá estar debidamente formado y certificado conforme a la legislación vigente en el momento de la intervención.


\end{document}